%!TEX program = uplatex
\RequirePackage{plautopatch}
\documentclass[uplatex,dvipdfmx]{jlreq}

\usepackage{amsmath,amssymb,amsthm}

% --- ここから本文 ---
\title{参考文献 (BibTeX) テンプレート}
\author{あなたのお名前}
\date{\today}

\begin{document}
\maketitle

\section{はじめに}
これは、別ファイル（\texttt{myref.bib}）を用いて参考文献リストを自動生成・管理するためのテンプレートです。

\section{引用テスト}
例えば、TeXの根幹たるコンセプトやアルゴリズムについては、開発者自身の著書 \cite{knuth84} が有名です。
また、日本における実践的な環境構築やマークアップの作法については、奥村・黒木による「美文書作成入門」 \cite{okumura2023} に詳しくまとまっています。

\vspace{2em}

% ==============================
% 参考文献リストの出力
% ==============================
% 論文で最も標準的な「出現順」のスタイル（junsrt）を指定
\bibliographystyle{junsrt}

% 読み込む .bib ファイルの名前を拡張子なしで指定
\bibliography{myref}

% ※注意: ファイルを分割・引用しているため、複数回コンパイルが必要です。
% 本ガイドの .latexmkrc を設定済みであれば、自動で全て実行・解決されます。
\end{document}
